\documentclass{article}
\usepackage{graphicx} % Required for inserting images
\usepackage[utf8]{inputenc}
\usepackage[russian]{babel}
\usepackage{verbatim}

\title{Документация проэкта}
\author{Цветков Владимир}

\begin{document}

\maketitle

\section{make\_tea(water\_ml, sugar\_spoons)}

\subsection{Описание}
Эта функция условно "заваривает чай" по вашему рецепту, принимая в себя количество миллилитров воды и ложек сахара для работы. Оба значения integer.

\subsection{Примеры вызова функции}

\begin{verbatim}
    make_tea(200, 2)
    # Output: Чай: 200 мл воды, 2 ложки сахара

    make_tea(100, 4)
    # Output: Чай: 100 мл воды, 4 ложки сахара
\end{verbatim}

\section{buy\_products(money, price\_per\_item, items\_count}
\subsection{Описание}
Функция принимает 3 значения: (int) Сумму ваших денег, (int) стоимость одного предмета и (int) сколько предметов нужно приобрести.
Возвращает функция одно из двух: Сдачу, оставшуюся после покупки, или предупреждение о том, что денег не хватает
\subsection{Примеры вызова функции}
\begin{verbatim}
    buy_products(100, 10, 5)
    # Output: 50

    buy_products(200, 40, 4)
    # Output: 80
\end{verbatim}

\section{check(day)}
\subsection{Описание функции}
Функция принимает (str) название нынешнего дня и проверяет, выходной он, или нет.
Возвращает True/False

\subsection{Примеры вызова функции}
\begin{verbatim}
    check("Суббота")
    # Output: True

    check("Понедельник")
    # Output: False
\end{verbatim}

\section{isEven(number)}
\subsection{Описание функции}
Функция принимает любое (int/float) число и проверяет, чётное оно, или нет.
Возвращает True/False
\subsection{Примеры вызова функции}
\begin{verbatim}
    isEven(2)
    # Output: True

    isEven(347)
    # Output: False
\end{verbatim}

\section{divByThree(start, end, step=1}
\subsection{Описание функции}
Функция принимает значения начала диапазона (int) start и его конца (int) end, также (int) step - шаг между каждым числом, всегда 1.
Возвращает список чисел, которые делятся на 3 в данном диапазоне.

\subsection{Примеры вызова функции}
\begin{verbatim}
    divByThree(0, 10, 1)
    # Output: [0, 3, 6, 9]

    divByThree(0, 2, 1)
    # Output: [0]
\end{verbatim}

\section{sentiment\_score(text, positive\_words, negative\_words)}
\subsection{Описание функции}
Функция принимает 3 значения: (str) сообщение, у которого нужно сосчитать счёт сентиментальности, (2 list with str) список позитивных и список отрицательных слов.
Выводит (int) счёт сентиментальности текста
\subsection{Примеры вызова функции}
\begin{verbatim}
    sentiment_score("I hate that pineapple!", [love, like, adore], [hate, disdain, pineapple, dislike])
    # Output: -3
\end{verbatim}

\section{binaryTransform(decimal\_number, counter=0}
\subsection{Описание функции}
Функция принимает 2 значения: (int) десятичное число, которое нужно трансформировать в двоичное и (int) счётчик, равный нулю.
Возвращает функция (str) бинарное значение десятичного числа

\subsection{Примеры вызова функции}
\begin{verbatim}
    binaryTransform(5, 0)
    # Output: '101'
\end{verbatim}

\section{История коммитов, хэши}
\begin{verbatim}
    0fec54 (HEAD -> doc_561991) Подправил функцию 3
    9949e63 Подправил последнюю функцию и добавил ей описание
    341e9da Сделал функцию счётчика сентиментальности текста попроще для прочтения, добавил описание
    6da1705 Поподробнее описал и изменил функцию 4, расписал переменные
    6a77790 Поподробнее описал и изменил функцию 3
    b959b94 Добавил описание ко второй функции и описал входную переменную поподробнее
    d66e71c Добавил описание функции №1 (о покупке товаров)
\end{verbatim}
\end{document}
